% GitHub Repo and Documentation: https://github.com/celiobjunior/resume-template
% Copyright © 2025 Celio B Junior. All rights reserved.
%
% Licensed under the Apache License, Version 2.0 (the "License");
% you may not use this file except in compliance with the License.
% You may obtain a copy of the License at
%
%     http://www.apache.org/licenses/LICENSE-2.0
%
% This template follows best practices from README.md

% Início do documento LaTeX: define tipo e formato do documento
% a4paper = tamanho A4, 10pt = tamanho base da fonte
\documentclass[a4paper,10pt]{article}

% --- PACOTES: BASE / IDIOMA ---
\usepackage[utf8]{inputenc}
\usepackage[T1]{fontenc}
\usepackage[portuguese]{babel}

% --- PACOTES: LAYOUT / TIPOGRAFIA ---
\usepackage{geometry}
\usepackage{parskip}
\usepackage{microtype}
\usepackage{enumitem}
\usepackage{titlesec}
\usepackage{array}
\usepackage{tabularx}

% --- PACOTES: LINKS / BOOKMARKS ---
\usepackage{hyperref}
\usepackage{bookmark}
\usepackage{xurl}

% --- CONFIGURAÇÃO: MARGENS / ESPAÇAMENTO ---
\geometry{top=1.5cm, bottom=1.5cm, left=1.5cm, right=1.5cm} % Margens da página
\setcounter{secnumdepth}{0}                                 % Remove numeração de seção
\setlist[itemize]{
    leftmargin=0.75em, % Recuo da lista (mais perto da margem da página)
    itemsep=0.2em,     % Espaço entre itens
    topsep=0.25em,     % Espaço antes/depois da lista
    parsep=0em,        % Espaço entre parágrafos no mesmo item
    partopsep=0em      % Espaço extra ao iniciar a lista
}

% Remove números de página e cabeçalhos para visual limpo do currículo
\pagestyle{empty}

% --- CONFIGURAÇÃO: METADADOS / LINKS ---
\hypersetup{
    pdftitle={CV Seu Nome},
    pdfauthor={Seu Nome},
    colorlinks=true,
    linkcolor=black,
    urlcolor=black,
    citecolor=black,
    bookmarksdepth=1 
}

% Desabilita numeração das sections
\setcounter{secnumdepth}{0}

% Formata os cabeçalhos de cada section e coloca uma linha em baixo
\titleformat{\section}
{\Large\bfseries}
{}
{0em}
{}
[\titlerule\vspace{0.5ex}]
\titleformat{\subsection}
{\normalsize\bfseries}
{}
{0em}
{}
\titlespacing*{\subsection}{0pt}{0.35em}{0.15em}

% Macro para entradas do currículo com bookmark no PDF
\newcounter{cventry}
\newcommand{\cventry}[4]{%
    \refstepcounter{cventry}%
    \phantomsection%
    \pdfbookmark[2]{#1}{cventry-\thecventry}%
    \noindent\begin{tabularx}{\textwidth}{@{}>{\raggedright\arraybackslash}X >{\raggedleft\arraybackslash}X@{}}
    \textbf{#1} & #2 \\
    \textit{#3} & \textit{#4} \\
    \end{tabularx}
}

% --- INÍCIO DO DOCUMENTO ---
\begin{document}

% --- CABEÇALHO ---
% Substitua com suas informações pessoais
\begin{center}
    {\LARGE \textbf{Raphael Augusto}} 
    \\ [0.1cm]
    Brasília, DF
    {\textbullet}
    Email: \href{mailto:raphael.augusto570@gmail.com}{raphael.augusto570@gmail.com} 
    {\textbullet}
    \href{https://www.linkedin.com/in/raphael-augusto-melo}{linkedin.com/in/raphael-augusto-melo} 
    {\textbullet}
    \href{https://github.com/usuario}{github.com/usuario}
\end{center}

% --- SEÇÕES ---

\section{Atividades de Liderança} 
    % Inclua trabalho voluntário, ligas acadêmicas ou cargos de liderança
    \subsection*{\texorpdfstring{
            \textbf{Nome da Empresa/Organização} \hfill Localização (ex: Remoto ou Estado, País)
        }{
            Nome da Empresa (Liderança) -- Localização
        }}
    \textit{Seu Cargo/Posição \hfill Data de Início - Data de Término (ex: Jan 2024 - Atual)}
        \begin{itemize}
            \item Desenvolveu sistema de monitoramento de tráfego humano em ambientes industriais utilizando YOLO e Python, implementando detecção em tempo real de pessoas em zonas de risco configuráveis por polígonos --- validado em vídeos de câmeras de segurança de ferrovias e de ambientes industriais em diferentes condições de iluminação.
            \item Desenvolveu pipeline completo de dados para gestão interna de colaboradores na arquitetura Medallion (Bronze, Silver e Gold) no Azure Databricks com PySpark, centralizando dados de funcionários para otimizar alocação em projetos e controle de disponibilidade por período de férias.
            \item Desenvolveu assistente de IA via WhatsApp para consultores de campo da Unilever Food Solutions (UFS), substituindo formulários manuais por conversa guiada com: identificação de operadores por CNPJ, análise de cardápios (foto, PDF ou texto) com matching vetorial de produtos Unilever, enriquecimento de dados via SerpAPI, planejamento de visita com histórico do cliente, registro estruturado de desfechos e agendamento da próxima visita.
            \item Integrou Azure OpenAI (Responses API), Azure Whisper (transcrição de áudio), Azure Vision OCR (extração de cardápios e documentos) e Azure Data Factory em pipelines de dados, compondo um ecossistema Azure end-to-end de ingestão e análise.
        \end{itemize}

% ------

\section{Experiência}
    % Liste as suas experiência de trabalho ou acadêmicas, mais recente primeiro
    \subsection*{\texorpdfstring{
            \textbf{Nome da Empresa/Organização} \hfill Localização (ex: Remoto ou Estado, País)
        }{
            Nome da Empresa -- Localização
        }}
    \textit{Cargo (ex: Engenheiro de Software, Analista de Dados) \hfill Data de Início - Data de Término}
        \begin{itemize} 
            % Foque em conquistas, não apenas na parte técnica - use o método STAR (Situação, Tarefa, Ação, Resultado)
            \item Alcancei [resultado específico] através de [ação tomada], resultando em [impacto quantificado].
            \item Desenvolvi/Construí/Criei [o que você fez] usando [tecnologias/métodos], resultando em [benefício].
            \item Colaborei com [tamanho/tipo da equipe] para [alcançar o quê], melhorando [métrica] em [porcentagem/quantidade].
        \end{itemize}

    \subsection*{\texorpdfstring{
        	\textbf{Nome da Empresa/Organização} \hfill Localização (ex: Remoto ou Estado, País)
        }{
            Nome da Empresa -- Localização
        }}
    \textit{Cargo (ex: Desenvolvedor Júnior, Assistente de Pesquisa) \hfill Data de Início - Data de Término}
        \begin{itemize}
            \item Desenvolveu middleware serverless utilizando Cloudflare Workers e TypeScript, garantindo escalabilidade e alta disponibilidade.
            \item Integrou APIs de fornecedores com Shopify via Webhooks, automatizando 100\% das vendas com provisionamento de eSIM via QRCode e envio de e-mails automáticos após a compra.
            \item Acompanhou deploys e ajustes pós-entrega via CI/CD, garantindo rápida evolução do sistema.
        \end{itemize}
        
    \subsection*{\texorpdfstring{
        	\textbf{Universidade/Instituição de Pesquisa} \hfill Localização (ex: Remoto ou Estado, País)
        }{
            Universidade/Instituição de Pesquisa -- Localização
        }}
    \textit{Função (ex: Pesquisador de Graduação, Monitor) \hfill Data de Início - Data de Término}
        \begin{itemize}
            \item Desenvolveu plataforma modular de orquestração de workflows com integração de IA generativa (Google Gemini) para criação dinâmica de fluxos a partir de linguagem natural.
            \item Atuou como líder técnico do projeto, definindo stack e arquitetura; projeto premiado como destaque de inovação no Centro Universitário de Brasília (CEUB).
        \end{itemize}

% ------

\section{Habilidades}
    % Mantenha esta seção concisa e use o máximo de palavras chaves
    % que fazem sentido para a vaga que você deseja.
    \begin{itemize}
        \item \textbf{Técnicas:} Python, JavaScript, TypeScript, SQL, PySpark, FastAPI, Django REST, Node.js, NestJS, PostgreSQL, Azure Databricks (Delta Lake), Azure OpenAI, Azure Data Factory, Docker, CI/CD, LLMs, RAG, Computer Vision (YOLO), Power BI, LangGraph, LangChain
        \item \textbf{Idiomas:} Português (Nativo), Inglês (Avançado)
    \end{itemize}

% ------

% !!!!!!!!!!!!!!!!!!!!!!
% Mude esta section para o começo, depois do header e antes
% das suas experiências se você está procurando sua primeira
% vaga ou algum estágio. Do contrário, deixe como está.
% !!!!!!!!!!!!!!!!!!!!!!

\section{Educação}
    % Formação mais recente primeiro
    \subsection*{\texorpdfstring{
            \textbf{Nome da Sua Universidade} \hfill Localização (ex: Remoto ou Estado, País)
        }{
            Nome da Sua Universidade (Educação) -- Localização
        }}
    \textit{Título do Seu Curso (ex: Bacharelado em Ciência da Computação) \hfill Data de Início - Data de Término (ex: Jan 2020 - Dez 2024)}

\end{document}